% ---------------------------------------------------------------------------
\begin{frame}{Armstrong Axioms --- Deriving New FDs}
  Named after William W.\ Armstrong (1974).  These inference rules let us
  \textbf{derive} every FD that logically follows from a given set $F$.

  \begin{columns}[t]
    \begin{column}{0.48\textwidth}
      \textbf{The Three Base Axioms}
      \begin{enumerate}
        \item \textbf{Reflexivity:}\\
              If $B \subseteq A$, then $A \to B$.\\
              \textit{A set determines its own subsets.}
        \item \textbf{Augmentation:}\\
              If $A \to B$, then $AC \to BC$.\\
              \textit{Adding attributes to both sides preserves the FD.}
        \item \textbf{Transitivity:}\\
              If $A \to B$ and $B \to C$, then $A \to C$.
      \end{enumerate}
    \end{column}
    \begin{column}{0.48\textwidth}
      \textbf{Derived Rules (shortcuts)}
      \begin{enumerate}
        \setcounter{enumi}{3}
        \item \textbf{Union:}\\
              $A \to B,\; A \to C \;\Rightarrow\; A \to BC$
        \item \textbf{Decomposition:}\\
              $A \to BC \;\Rightarrow\; A \to B$ and $A \to C$
        \item \textbf{Pseudotransitivity:}\\
              $A \to B,\; CB \to D \;\Rightarrow\; CA \to D$
      \end{enumerate}
    \end{column}
  \end{columns}

  \vspace{0.5em}
  \begin{alertblock}{Completeness and Soundness}
    The Armstrong axioms are both \textbf{complete} (every FD that holds can be
    derived) and \textbf{sound} (every derived FD actually holds). They form
    the foundation for all normalisation algorithms.
  \end{alertblock}
\end{frame}

% ---------------------------------------------------------------------------
\begin{frame}{Armstrong Axioms --- Applied to the Music Example}
  Given the following FDs for our music catalogue:
  \begin{itemize}
    \item $\texttt{AlbumID} \to \texttt{ArtistID}$
    \item $\texttt{ArtistID} \to \texttt{DebutYear}$
    \item $\texttt{AlbumID} \to \texttt{Genre}$
    \item $\texttt{AlbumID} \to \texttt{ReleaseYear}$
  \end{itemize}

  \vspace{6pt}
  \begin{examples}{Deriving Further FDs}
    \begin{tabular}{lll}
      \toprule
      \textbf{Derived FD} & \textbf{Rule Used} & \textbf{From} \\
      \midrule
      $\texttt{AlbumID} \to \texttt{DebutYear}$ & Transitivity & FD 1 + FD 2 \\
      $\texttt{AlbumID} \to \texttt{Genre, ReleaseYear}$ & Union & FD 3 + FD 4 \\
      $\texttt{AlbumID} \to \texttt{ArtistID, Genre}$ & Union & FD 1 + FD 3 \\
      $\texttt{AlbumID, TrackNo} \to \texttt{ArtistID}$ & Augmentation & FD 1 \\
      \bottomrule
    \end{tabular}
  \end{examples}
\end{frame}

% ---------------------------------------------------------------------------
\begin{frame}{Attribute Closure $A^+$}
  \begin{block}{Definition}
    The \textbf{closure} of a set of attributes $A$ under FDs $F$, written
    $A^+$, is the set of \emph{all attributes} that $A$ functionally determines
    (directly or transitively).
  \end{block}

  \begin{block}{Algorithm: Compute $A^+$}
    \begin{enumerate}
      \item Start: $A^+ := A$
      \item Repeat until no change:\\
            for each FD $X \to Y$ in $F$: if $X \subseteq A^+$, add $Y$ to $A^+$
    \end{enumerate}
    If $A^+ = R$, then $A$ is a \textbf{superkey}.
  \end{block}

  \begin{examples}{Example}
    $R = \{A,B,C,D,E\}$, \quad $F = \{A \to BC,\; B \to D,\; D \to E\}$

    Compute $\{A\}^+$:
    \begin{tabular}{ll}
      Start: & $\{A\}$ \\
      $A \to BC$: & $\{A,B,C\}$ \\
      $B \to D$: & $\{A,B,C,D\}$ \\
      $D \to E$: & $\{A,B,C,D,E\} = R$ \\
    \end{tabular}

    $\{A\}^+ = R \;\Rightarrow$ \textbf{A is a superkey!}
  \end{examples}
\end{frame}

% ---------------------------------------------------------------------------
\begin{frame}{Keys --- Formal Definitions}
  \begin{itemize}
    \item \textbf{Superkey:} A set $A \subseteq R$ such that $A^+ = R$.
          May contain redundant attributes.
    \item \textbf{Candidate Key:} A \emph{minimal} superkey --- removing any
          single attribute makes it no longer a superkey.
    \item \textbf{Primary Key:} The candidate key chosen as the principal identifier.
    \item \textbf{Prime Attribute:} An attribute that is part of \emph{at least one} candidate key.
    \item \textbf{Non-Prime Attribute:} An attribute \emph{not} part of any candidate key.
          These are the focus of most normalisation rules.
  \end{itemize}

  \begin{block}{Hierarchy}
    \[ \text{Superkeys} \;\supset\; \text{Candidate Keys} \;\supset\; \{\text{Primary Key}\} \]
  \end{block}

  \begin{alertblock}{How to Find All Candidate Keys}
    For each subset $A$ of $R$: if $A^+ = R$ and no proper subset of $A$ has
    the same property, then $A$ is a candidate key.
    Start with single attributes, then pairs, triples, \ldots
  \end{alertblock}
\end{frame}

% ---------------------------------------------------------------------------
\begin{frame}{Exercise: Closure and Candidate Keys}
  \begin{block}{Task}
    Given $R = \{A, B, C, D\}$ and $F = \{A \to B,\; B \to C,\; CD \to A\}$:
    \begin{enumerate}
      \item Compute $\{A\}^+$, $\{B\}^+$, $\{D\}^+$, $\{AD\}^+$, $\{BD\}^+$, $\{CD\}^+$
      \item Identify all candidate keys of $R$
      \item Which attributes are prime? Which are non-prime?
    \end{enumerate}
  \end{block}

  \begin{alertblock}{Time: 10 minutes}
    Then we discuss the solutions together.
  \end{alertblock}
\end{frame}
